\documentclass[11pt]{article}
\usepackage[margin=1in]{geometry}
\usepackage{amsmath,amssymb,amsthm,booktabs,longtable,array,calc}
\usepackage[hidelinks]{hyperref}
\usepackage[T1]{fontenc}
\usepackage[utf8]{inputenc}
\newtheorem{theorem}{Theorem}[section]
\newtheorem{proposition}[theorem]{Proposition}
\newtheorem{lemma}[theorem]{Lemma}
\newtheorem{corollary}[theorem]{Corollary}
\theoremstyle{definition}
\newtheorem{definition}[theorem]{Definition}
\theoremstyle{remark}
\newtheorem*{remark}{Remark}
\providecommand{\tightlist}{\setlength{\itemsep}{0pt}\setlength{\parskip}{0pt}}
\begin{document}
\appendix
\section{Full Validity of the Canonical Pleat Orders}\label{appendix-c.-full-validity-of-the-canonical-pleat-orders}

The Hamiltonian arguments in Theorem 5.1 prove uniqueness among
precedence-compatible orders. This appendix proves existence by checking
every Butterfly family at arbitrary separation.

\subsection{\texorpdfstring{C.1. The order
\(\Omega_+(k)\)}{C.1. The order \textbackslash Omega\_+(k)}}\label{c.1.-the-order-omega_k}

For \(1\le q\le k\) the positions are \[
\operatorname{pos}(A_{2q-1})=4q-3,\quad
\operatorname{pos}(B_{2q-1})=4q-2,
\] \[
\operatorname{pos}(B_{2q})=4q-1,\quad
\operatorname{pos}(A_{2q})=4q.
\] All primitive M/V precedences point forward in this order. The sorted
endpoint intervals are \[
I(H_{2q-1})=[4q-3,4q-2],\qquad I(H_{2q})=[4q-1,4q],
\] \[
I(U_{2q-1})=[4q-3,4q],\qquad I(U_{2q})=[4q,4q+1],
\] \[
I(L_{2q-1})=[4q-2,4q-1],\qquad I(L_{2q})=[4q-1,4q+2],
\] where the even-seam formulas run only to \(q=k-1\).

For \(H/H\), distinct column intervals are strictly disjoint. For
same-parity \(U/U\) and \(L/L\), if \(q<r\) the right endpoint of the
\(q\) interval is strictly before the left endpoint of the \(r\)
interval; hence these pairs are disjoint.

For the cross-row \(U/L\) class there is no distance cutoff. For odd
seams \(i=2q-1\), \(j=2r-1\), \[
I(U_i)=[4q-3,4q],\qquad I(L_j)=[4r-2,4r-1].
\] If \(r=q\), the lower interval is strictly nested in the upper one;
if \(r<q\) it is strictly before it; if \(r>q\) it is strictly after it.
For even seams \(i=2q\), \(j=2r\), \[
I(U_i)=[4q,4q+1],\qquad I(L_j)=[4r-1,4r+2],
\] and the same trichotomy holds, with the upper interval nested in the
lower one when \(r=q\). Thus every same-target pair is nested or
disjoint for arbitrary \(q,r\).

\subsection{\texorpdfstring{C.2. The order
\(\Omega_-(k)\)}{C.2. The order \textbackslash Omega\_-(k)}}\label{c.2.-the-order-omega_-k}

Here \[
\operatorname{pos}(A_c)=c,\qquad
\operatorname{pos}(B_c)=2n-c+1.
\] The primitive directed relations are \(A_i\prec A_{i+1}\),
\(B_{i+1}\prec B_i\), and \(A_c\prec B_c\), all of which point forward.

For horizontal creases, \[
I(H_c)=[c,2n-c+1].
\] If \(c<d\), then \[
c<d<2n-d+1<2n-c+1,
\] so \(I(H_d)\) is strictly nested in \(I(H_c)\).

Upper vertical intervals are \(I(U_i)=[i,i+1]\). Distinct same-target
seams have equal parity and hence differ by at least two, so these
intervals are disjoint. Lower vertical intervals are \[
I(L_i)=[2n-i,2n-i+1]
\] and are likewise disjoint within a target parity.

Finally, \[
\max I(U_i)\le n<n+1\le\min I(L_j)
\] for all \(i,j\), so every cross-row \(U/L\) pair is separated, even
before target parity is imposed.

\subsection{C.3. Boundary and strictness
check}\label{c.3.-boundary-and-strictness-check}

The Butterfly predicate uses strict alternation. Genuine tested pairs
have disjoint wing sets, so their four endpoint positions are distinct.
At the left and right boundaries all displayed interval formulas stay
within \(1,\ldots,4k\); even seam formulas stop at \(n-2\). Therefore no
endpoint degeneration or omitted long-range case occurs. The families
\(H/H\), \(U/U\), \(L/L\), and \(U/L\) exhaust all same-target disjoint
crease pairs because horizontal target \(S\) is distinct from vertical
targets \(E,W\).

Hence both canonical pleat orders satisfy every global semantic
condition for every \(k\ge2\).

\end{document}
